\newpage
\section{LITERATURE REVIEW}

Wireless communication has greatly evolved since the experiments conducted by Marconi at the beginning of the twentieth century and today forms one of the integral components of modern communication systems~\cite{ref3}. Progress in the field of wireless communication networks has been driven primarily by the need for greater mobility, increased transmission rate, and efficient utilization of available bandwidth~\cite{ref4}. One of the crucial elements responsible for the performance of any wireless communication system is its antennas, whose design plays an important role in defining the range, data rate, and susceptibility to interference~\cite{ref5}. Today there are multiple wireless technologies in use, ranging from Wireless Fidelity (WiFi, IEEE 802.11) and Bluetooth, to Near Field Communication (NFC), Long-Term Evolution (LTE), and Fifth Generation (5G) wireless technology. There are also emerging wireless networks like LoRa that specialize in long-range communication, and each of these networks utilizes different ranges, data rates, and energy requirements. At the same time, the challenge of ensuring reliable long-range communication still persists because of various factors that may cause communication interference.

There are two main types of antennas: omnidirectional and directional. While the former offers wide coverage area, their main disadvantages include low gain, high interference, and less energy-efficient design. Directional antennas are able to emit the transmitted signal in specific directions, achieving increased gain, range, low interference, and energy efficiency~\cite{ref6}. As a result, they are extensively used in various long-range networks, for example in Wireless Sensor Networks (WSNs), Unmanned Aerial Vehicle (UAV) communication systems, Internet of Things (IoT), and Mobile Ad hoc Networks (MANETs). However, their effective use depends on precise alignment of transmitters and receivers with one another, requiring accurate orientation.

There exist several traditional methods for obtaining the optimal antenna orientation, including global positioning system (GPS)-based coordinate alignment and the signal tracking control method (TCM)~\cite{ref7,ref8}. While these methods are able to align the antennas, they involve additional hardware in the form of GPS modules, complicating the system structure. In addition, GPS-based methods are susceptible to signal fluctuations, making them unattractive for use indoors. Control theory-based approaches, such as imitation human intelligence control~\cite{ref9} and servo tracking, are also able to perform automatic antenna orientation~\cite{ref10} . However, they depend heavily on the accurate modeling of a specific system, making their implementation challenging.

In situations where position data is missing or unreliable, antenna alignment should be achieved through feedback from the signal, mainly RSSI. This is because RSSI-based alignment allows controlling the orientation of the node without the requirement to have information about the positions, thus reducing the costs associated with equipment~\cite{ref11,ref12}. However, early approaches based on RSSI feedback were based on exhaustive scanning of different orientations until the one with maximum RSSI was found. As can be seen from Table I, this approach is simple and inexpensive in terms of hardware required; however, it requires a lot of computations and, therefore, time to converge, making it impractical in real-time situations.

Other approaches to aligning directional antennas using signal feedback include estimation and control methods, such as Kalman filters, unscented Kalman filters, and fuzzy logic~\cite{ref13}. The use of estimation techniques improves robustness by accounting for fluctuations in the signal; however, such algorithms are complex and usually need additional information about the environment, either acquired with sensors or assumed before deployment. Therefore, additional hardware becomes necessary, decreasing the embeddability of the solution.

There are also autonomous systems that are able to orient themselves based on received signals (RSSI). They show the potential of a simple and inexpensive mechanical mechanism to adjust the orientation based on signal feedback~\cite{ref14}.. However, due to heuristics involved in such mechanisms, they cannot cope well with noisy and dynamically changing environment, besides being slow in their operation.

One approach to RSSI-based alignment is based on a nonlinear distributed control law derived from RSSI feedback for GPS-free antenna alignment in mobile platforms \cite{ref11}. In this algorithm, the orientation is estimated through perturbations of antenna position to find the direction of maximum signal intensity. However, such an approach is computationally costly, since it requires perturbation estimation and introduces sensitivity to noise, which makes the real-time implementation impractical for embedded hardware.

One recent trend in addressing the problem with RSSI-based antenna alignment is using machine learning (ML). Supervised learning methods using MLP, CNN, and LSTM networks can achieve very good results in predicting RSSI in different orientations and estimating the optimal orientation~\cite{ref15}. However, they require considerable amounts of data, offline training, and precise measurement of data. In turn, all this imposes considerable costs associated with deploying the system in practice.

Along the same lines, DRL approaches, including Deep Q-Networks (DQN), have been utilized for beam alignment in millimeter-wave (mmWave) and 5G communications ~\cite{ref17}. Even though such techniques work efficiently in complex and dynamic settings, they rely on the use of neural network algorithms that entail significant memory consumption and computational effort. This exceeds the capacities of microcontroller-level computing, making DRL-based approaches inapplicable in the context of embedded platforms lacking hardware acceleration capabilities.

Also, the use of MARL approaches in optimizing antenna tilt angle in cellular networks has been proposed~\cite{ref18}. Even though MARL can optimize parameters at various nodes simultaneously, it requires network-wide coordination and extensive computational power, thus increasing complexity and restricting the usage scenarios of this method to standalone embedded platforms.

Overall, there are clear tradeoffs between the effectiveness and practicality of current solutions. Control-theoretic and estimation-based approaches show good performance theoretically; however, they require precise models and higher computational complexity. On the other hand, exhaustive scanning is a simple yet inexpensive process but cannot ensure efficient real-time operation. Supervised and deep learning-based algorithms provide better accuracy but entail data storage costs and significant computational requirements. Moreover, many existing solutions emphasize antenna arrays and beamforming algorithms, which increase both hardware cost and complexity in comparison to mechanically steered antennas.

\newpage
\setlength{\LTpre}{1cm}   % no space before table
\setlength{\LTpost}{1cm}  % no space after table
\renewcommand{\arraystretch}{1.4}
\setlength{\tabcolsep}{8pt}

\begin{longtable}{|
    >{\raggedright\arraybackslash}p{4.3cm} |
    >{\raggedright\arraybackslash}p{4.3cm} |
    >{\raggedright\arraybackslash}p{4.1cm} |}

\caption{Summary of Literature Review}
\label{tab:literature-review} \\
\hline
\multicolumn{1}{|c|}{\parbox[c][2em][c]{4cm}{\centering \textbf{Title}}} &
\multicolumn{1}{c|}{\parbox[c][2em][c]{4.5cm}{\centering \textbf{Methodology}}} &
\multicolumn{1}{c|}{\parbox[c][2em][c]{4.5cm}{\centering \textbf{Key Inferences}}} \\
\hline

\textit{RSSI-Based Distributed Control to Align Directional Antenna Pairs for UAV Communication} (2024) \newline
Y. Wan et al.
&
Control-theoretic approach using nonlinear static state feedback. RSSI partial derivatives estimated via antenna motion.
&
Demonstrates GPS-free RSSI-based antenna alignment with provable convergence. Strong theory, but requires accurate derivative estimation, difficult on low-cost embedded hardware.
\\ \hline

\textit{Received Signal Strength Indicator-Based Decentralised Control for Robust Long-Range Aerial Networking Using Directional Antennas} (2017) \newline
Yan et al.
&
Fusion of RSSI and GPS using Unscented Kalman Filter and fuzzy logic.
&
Improves robustness but depends on additional sensors and antenna models. Not suitable for minimal hardware systems.
\\ \hline

\textit{Self-Orientation of Directional Antennas Assisted by Mobile Robots} (2012) \newline
Min et al.
&
Pattern-based RSSI scanning using pan-tilt antenna mounts.
&
Validates RSSI-only alignment, but exhaustive scanning is slow and unsuitable for real-time systems.
\\ \hline

\textit{SmartAntenna: Enhancing Wireless Range with Autonomous Orientation} (2024) \newline
Swann et al.
&
Servo-mounted antenna using RSSI feedback and heuristic search.
&
Demonstrates low-cost feasibility, but performance degrades in noisy environments and lacks adaptive learning.
\\ \hline

\textit{Comparative Analysis of Machine Learning Algorithms for Antenna Alignments} (2024) \newline
Shakya et al.
&
Supervised ML (MLP, CNN, LSTM) trained on S-parameters.
&
High accuracy in controlled environments, but requires large datasets and offline training.
\\ \hline

\textit{Learning-Based Beam Alignment for Uplink mmWave UAVs} (2023) \newline
Susarla et al.
&
Deep Q-Network for beam alignment using antenna arrays.
&
Effective in mmWave systems, but unsuitable for single-antenna, low-cost embedded platforms.
\\ \hline

\end{longtable}


The following research highlights a significant research gap related to the design of an efficient antenna alignment system driven by RSSI and having the capability to work in real-time on embedded systems. The solution provided for addressing the aforementioned problem takes into consideration the use of minimum hardware, RSSI only, and the application of reinforcement learning techniques. Compared to other techniques, this technique does not require GPS, models, big data, and high computational power, thus being much more realistic in its trade-offs.